\documentclass[10pt,a4paper]{amsart}
\usepackage[margin=19mm]{geometry}
\usepackage{amsmath,amssymb,mathtools}
\usepackage[scr=rsfs,cal=euler]{mathalfa}
\usepackage{xfrac}
\usepackage[unicode,hidelinks,bookmarks=false]{hyperref}
\newcommand{\ie}{i.e.\ }
\newcommand{\cf}{cf.\ }
\newcommand{\resp}{respectively\ }
\newcommand{\Subsection}{\subsection}
\providecommand{\nobreakdash}{\nobreak\hbox{-}}
\setlength{\emergencystretch}{2em}
\multlinegap0pt
\usepackage{bm}
\usepackage{amssymb}
\usepackage{xcolor}
\usepackage{mathtools}
\usepackage{algorithm}\usepackage{algpseudocode}
\usepackage{tikz}
\newtheorem{lemma}{Lemma}
\newtheorem{proposition}{Proposition}
\title{Model order reduction for parametrized variational inequalities: application to crowd motion}
\author{Giulia Sambataro, Virginie Ehrlacher}
\date{Source: arXiv:2605.04037v1 (CC BY 4.0)}
\begin{document}
\maketitle

\noindent\emph{Source and licence.} Model order reduction for parametrized variational inequalities: application to crowd motion. Version arXiv:2605.04037v1 (CC BY 4.0), distributed by arXiv under the Creative Commons Attribution 4.0 International licence, http://creativecommons.org/licenses/by/4.0/, as recorded in the arXiv licence metadata for this exact version and shown on the abstract page https://arxiv.org/abs/2605.04037v1. Full licence notice: \texttt{licenses/arxiv-2605.04037-CC-BY-4.0.txt}.\par\smallskip

\noindent\textit{Editorial scope.} Counted unit: the article's proposition prop:lambdas-explicit (source0.tex lines 1376--1399). The article's complete original proof of it is delivered with it (source0.tex lines 1401--1479, one \texttt{proof} environment of 79 lines). No mathematics is written, repaired or re-derived: the statement and the proof are the article's own lines, in the article's own notation, and no retained line is substituted or re-flowed.  Retained uncounted as the source-local context the proof needs: lines 1270--1274; lines 1353--1365.  Nothing was omitted from any retained span, so the package carries no omission record.  No retained line contains a citation, so the wrapper carries no bibliography and no bibliography entry.  No retained line contains a figure environment, an image-inclusion command or drawing code, so no image is delivered and the package declares no asset dependency.  Editorial formatting only, in addition: an AMS \texttt{amsart} preamble that restates the article's own declarations and macros used by the retained lines, namely the environments \texttt{lemma}, \texttt{proposition} on separate counters, together with the standard packages those lines need.  The retained environments print in this document as Lemma 1, Proposition 1 in order of appearance, whereas the article prints the same items with the numbering of its own full document.  The complete native source of the article as distributed in the arXiv e-print for this exact version is bundled unchanged as \texttt{sources/199833/source0.tex}.
\begin{equation}\label{eq:minimisation}
\min_{(\lambda_1,\lambda_2)\in\mathbb{R}_+^2}
\Bigl\|p\,\varphi_a-\lambda_1\,\overline{p}\,\varphi_{\overline{a}}
-\lambda_2\,\overline{p}\,\varphi_{\underline{a}}\Bigr\|_{L^2(\mathbb{R})}^2,
\end{equation}

\begin{lemma}\label{lem:gram}
The Gram matrix
\[
\mathbf{G}:=\overline{p}^2\begin{pmatrix}
\mathcal{A} & \mathcal{C}\\
\mathcal{C} & \mathcal{B}
\end{pmatrix}
\]
is positive definite. In particular,
\[
\mathcal{A}\mathcal{B}-\mathcal{C}^2>0.
\]
\end{lemma}

\begin{proposition}\label{prop:lambdas-explicit}
For every $(a,p)\in[\underline{a},\overline{a}]\times(0,\infty)$, the minimizer of
problem \eqref{eq:minimisation} is given by
\begin{equation}\label{eq:lambda1-formula}
\lambda_1(a,p)
=
\frac{p}{\overline{p}}\,\frac{\mathcal{B}\,\mathcal{X}(a)-\mathcal{C}\,\mathcal{Y}(a)}
                             {\mathcal{A}\mathcal{B}-\mathcal{C}^2},
\end{equation}
\begin{equation}\label{eq:lambda2-formula}
\lambda_2(a,p)
=
\frac{p}{\overline{p}}\,\frac{\mathcal{A}\,\mathcal{Y}(a)-\mathcal{C}\,\mathcal{X}(a)}
                             {\mathcal{A}\mathcal{B}-\mathcal{C}^2}.
\end{equation}
In particular, at the endpoints of the parameter domain:
\[
\lambda_1(\overline{a},p)=\frac{p}{\overline{p}},\quad
\lambda_2(\overline{a},p)=0,
\qquad
\lambda_1(\underline{a},p)=0,\quad
\lambda_2(\underline{a},p)=\frac{p}{\overline{p}}.
\]
\end{proposition}

\begin{proof}
The optimality conditions for problem \eqref{eq:minimisation} (ignoring the
non-negativity constraints, which will be verified a posteriori) read:
\[
\mathbf{G}\begin{pmatrix}\lambda_1\\ \lambda_2\end{pmatrix}
=
\begin{pmatrix}
\langle\lambda_{a,p},\psi_1\rangle_{L^2}\\
\langle\lambda_{a,p},\psi_2\rangle_{L^2}
\end{pmatrix}.
\]
By the separability $\lambda_{a,p}=p\,\varphi_a$ and $\psi_i=\overline{p}\,\varphi_{\cdot}$,
the right-hand sides are
\[
\langle\lambda_{a,p},\psi_1\rangle_{L^2}=p\overline{p}\,\mathcal{X}(a),
\qquad
\langle\lambda_{a,p},\psi_2\rangle_{L^2}=p\overline{p}\,\mathcal{Y}(a).
\]
The system becomes
\[
\overline{p}^2
\begin{pmatrix}\mathcal{A} & \mathcal{C}\\ \mathcal{C} & \mathcal{B}\end{pmatrix}
\begin{pmatrix}\lambda_1\\ \lambda_2\end{pmatrix}
=
p\overline{p}
\begin{pmatrix}\mathcal{X}(a)\\ \mathcal{Y}(a)\end{pmatrix}.
\]
By Lemma~\ref{lem:gram} the matrix is invertible; dividing by $\overline{p}^2$
and inverting gives
\[
\begin{pmatrix}\lambda_1\\ \lambda_2\end{pmatrix}
=
\frac{p/\overline{p}}{\mathcal{A}\mathcal{B}-\mathcal{C}^2}
\begin{pmatrix}
\mathcal{B} & -\mathcal{C}\\
-\mathcal{C} & \mathcal{A}
\end{pmatrix}
\begin{pmatrix}\mathcal{X}(a)\\ \mathcal{Y}(a)\end{pmatrix},
\]
which yields \eqref{eq:lambda1-formula} and \eqref{eq:lambda2-formula}.

\medskip\noindent
\emph{Endpoint cases.}
For $a=\overline{a}$: $\mathcal{X}(\overline{a})=J(\overline{a},\overline{a})=\mathcal{A}$
and $\mathcal{Y}(\overline{a})=J(\underline{a},\overline{a})=\mathcal{C}$, so
\[
\lambda_1(\overline{a},p)
=\frac{p}{\overline{p}}\frac{\mathcal{B}\mathcal{A}-\mathcal{C}^2}
                            {\mathcal{A}\mathcal{B}-\mathcal{C}^2}
=\frac{p}{\overline{p}},
\qquad
\lambda_2(\overline{a},p)
=\frac{p}{\overline{p}}\frac{\mathcal{A}\mathcal{C}-\mathcal{C}\mathcal{A}}
                            {\mathcal{A}\mathcal{B}-\mathcal{C}^2}=0.
\]
For $a=\underline{a}$: $\mathcal{X}(\underline{a})=J(\underline{a},\overline{a})=\mathcal{C}$
and $\mathcal{Y}(\underline{a})=J(\underline{a},\underline{a})=\mathcal{B}$, so
\[
\lambda_1(\underline{a},p)
=\frac{p}{\overline{p}}\frac{\mathcal{B}\mathcal{C}-\mathcal{C}\mathcal{B}}
                            {\mathcal{A}\mathcal{B}-\mathcal{C}^2}=0,
\qquad
\lambda_2(\underline{a},p)
=\frac{p}{\overline{p}}\frac{\mathcal{A}\mathcal{B}-\mathcal{C}^2}
                            {\mathcal{A}\mathcal{B}-\mathcal{C}^2}=\frac{p}{\overline{p}}.
\]
In particular, for $p=\overline{p}$ we recover $\lambda_1(\overline{a},\overline{p})=1$
and $\lambda_2(\underline{a},\overline{p})=1$, consistent with the greedy
selections $\psi_1=\lambda_{\overline{a},\overline{p}}$ and
$\psi_2=\lambda_{\underline{a},\overline{p}}$.

\medskip\noindent
Finally, since $\mathcal{X}$ and $\mathcal{Y}$ are strictly increasing (established
in the next section), formulas \eqref{eq:lambda1-formula}--\eqref{eq:lambda2-formula}
yield non-negative coefficients for all $a\in[\underline{a},\overline{a}]$; hence
the minimizer of the constrained problem over $\mathbb{R}_+^2$ coincides with the
unconstrained solution above.
\qed
\end{proof}

\end{document}
